\section{导子和代数同态的关系}

记号约定:$A$是分次$K$-代数,$E$是分次$K$-模,配备了$A$的左乘和右乘诱导的作用,$\delta\in\Delta$,$\epsilon$是交换因子.

\begin{definition}{}{}
	定义$K$-模$A\boxplus E\left(\delta\right)$上的分次$K$-代数结构如下:对$A$的每个齐次元$a$和每个$a^{\prime}\in A,x,x^{\prime}\in E\left(\delta\right)$,
	\begin{align*}
		\left(a,x\right)\cdot\left(a^{\prime},x^{\prime}\right)\coloneq\left(aa^{\prime},xa^{\prime}+\epsilon\left(\delta,\deg\left(a\right)\right)ax^{\prime}\right).
	\end{align*}
	我们称典范投影$p:A\boxplus E\left(\delta\right)\to A,\left(a,x\right)\mapsto a$是$A\boxplus E\left(\delta\right)$的增广.
\end{definition}

\begin{remark}{}{}
	如果$g:A\to A\boxplus E\left(\delta\right)$是$p$的左逆,则存在$\delta$次分次$K$-线性映射$f:A\to E$满足$\forall x\in A\left(g\left(x\right)=\left(x,f\left(x\right)\right)\right)$.
\end{remark}

\begin{proposition}{导子和代数同态的等价性}{}
	设$f:A\to E$是$\delta$次分次$K$-线性映射,则$f$是$\epsilon$-导子$\iff$ $A\to A\boxplus E\left(\delta\right),x\mapsto\left(x,f\left(x\right)\right)$是分次$K$-代数同态.
\end{proposition}

\begin{proposition}{}{}
	$A\boxplus E\left(\delta\right)$是含幺结合代数的等价条件是:
	\begin{enumerate}
		\item $A$是含幺结合代数,
		\item 映射$A\times E\to E,\left(a,x\right)\mapsto ax$和$E\times A\to A,\left(x,a\right)\mapsto xa$使得$E$是典范$\left(A,A\right)$-双模.
	\end{enumerate}
	这两个条件成立时,$A\boxplus E\left(\delta\right)$的单位元是$\left(1_{A},0\right)$.
\end{proposition}



















































